\scanpage{135}

Théorème de \handwritten{- - - - - -} (rapporté par Husemöller).

Soit $A$ une bialgèbre sur le corps $k$, $A$ associative et unitaire pour ses deux structures, ``cocomplète'' pour sa filtration canonique de comodule (duale de la filtration par les puissances de l'idéal d'augmentation du dual; de façon précise, écrivant $\overset{n}{\otimes}A=\overset{n}{\otimes}A^{\handwritten{+}}+$ fo\deleted{r}u\uline{r}bis, d'où via la m.ème itéré\deleted{\unsure{e|r}} du morphisme diagonal, une application $A\handwritten{\longrightarrow}\overset{n}{\otimes}I^{\handwritten{+}}$, dont le noyau\inserted{\handwritten{\unsure{intersect.|intersect°} avec $A^{+}$}}\figure{\begin{tikzpicture}[baseline=-.6ex,x=1cm,y=1cm]
\draw (.06,-.24) .. controls (.02,.12) and (.34,.18) .. (.75,.2) -- (3.8,.2) .. controls (4.2,.22) and (4.2,.32) .. (4.16,.43);
\end{tikzpicture}} est $\mathrm{Filt}_{m}(A)$).\margin{\handwritten{cela entraine qu'il y a un antipodisme ?}} (NB l'hypothèse cocomplète est satisfaite pour une coalgèbre \inserted{avec \unsure{unité|unite}} à degrés $\handwritten{\geq}$ $0$ réduite à $k$ en degré\deleted{s} $0$, associative et unitaire, - parait-il !). Soit $M$ un $A$-module, muni d'un homomorphisme diagonal \inserted{sur $k$}\figure{\begin{tikzpicture}[baseline=-.6ex,x=1cm,y=1cm]
\draw (.06,-.24) .. controls (.02,.12) and (.34,.18) .. (.75,.2) -- (1.2,.2) .. controls (1.6,.22) and (1.6,.32) .. (1.56,.43);
\end{tikzpicture}} unitaire (peut-être pas associatif), $M\handwritten{\longrightarrow}M\otimes_{k}M$, qui soit $A$-linéaire,\margin{\handwritten{\unsure{6|b} faut il mett. dans les hyp. que $A\otimes M\longrightarrow M$ \unsure{compatible|comp.tible} aux \unsure{counités|unités} !}\figure{\begin{tikzpicture}[x=1cm,y=1cm]\draw (0,3.2) .. controls (.9,3.35) and (2.3,3.3) .. (2.35,2.95) .. controls (2.25,1.1) and (1.8,-.1) .. (0,0);\end{tikzpicture}}}\inserted{\handwritten{(pour la structure de $A$-module sur $M\unsure{\otimes|\ominus}M$ provenant de $A\longrightarrow A\unsure{\otimes|\ominus}A$; il revient au \unsure{même|m.} de dire que $A\unsure{\otimes|\ominus}M\longrightarrow M$ comp. avec \unsure{applications|applictions|applic\textsuperscript{tions}} diagonales)}} Supposons de plus donné une ``unité'' i.e. un homomorphisme $k\handwritten{\longrightarrow}M$ qui soit un homomorphisme de coalgèbres unitaires, i.e. $\underset{M}{\mathrm{l}}$'image \inserted{$e_M$}\figure{\begin{tikzpicture}[baseline=-.6ex,x=1cm,y=1cm]
\draw (.13,-.2) .. controls (.46,.2) and (.94,.26) .. (1.01,.67);
\end{tikzpicture}} de $1$ satisfait $\handwritten{\epsilon}_{M}(e_{M})=1$, $\handwritten{\Delta}_{M}(e_{M})=e_{M}\otimes e_{M}$. Supposons également $M$ cocomplète, et qu\deleted{\unsure{'|l}}e\deleted{on} l'application $u\handwritten{\mapsto}u.e$ de $A$ dans $M$ soit injective. Alors $M$ est un $A$-module libre, \inserted{\handwritten{ayant une base dont fait partie $e$;}}\figure{\begin{tikzpicture}[baseline=-.6ex,x=1cm,y=1cm]
\draw (.06,-.24) .. controls (.02,.12) and (.34,.18) .. (.75,.2) -- (7.0,.2) .. controls (7.4,.22) and (7.4,.32) .. (7.36,.43);
\end{tikzpicture}} de façon plus précise, posant $\overset{\handwritten{-}}{M}=M\otimes_{A}k$, il existe un isomorphisme (non canonique ...) de $A$-modules $M\handwritten{\unsure{\deleted{\longrightarrow}\overset{\sim}{\longleftarrow}|\deleted{=}\overset{\sim}{\longleftarrow}|\overset{\sim}{\longleftrightarrow}}}A\otimes_{k}\overset{\handwritten{-}}{M}$\handwritten{, tel que le diagramme}
\margin{\handwritten{?}}
\handwritten{\diagram{\begin{tikzcd}[ampersand replacement=\&,column sep=small,row sep=small]
M \arrow[dr] \& \unsure{\deleted{M}A|A|\mathcal{A}}\otimes_{k}\overline{M} \arrow[l,"\sim"'] \arrow[d] \\
\& \overline{M}
\end{tikzcd}}}
\handwritten{commute, et que \unsure{...|. -}}

Exemple \handwritten{1}. Supposons l'application diagonale de $A$ commutative. Alors $A$ peut être considéré comme l'hyperalgèbre d'un groupe formel connexe \inserted{$G$}\figure{\begin{tikzpicture}[baseline=-.6ex,x=1cm,y=1cm]
\draw (-.08,.42) .. controls (-.35,.25) and (.65,.47) .. (.58,-.24) -- (1.02,.65);
\end{tikzpicture}} sur $k$. Si $\handwritten{\Delta}_{M}$ est associatif et commutatif, $M$ aussi est formé des distributions à l'origine d'une variété formelle connexe \inserted{$V$}\figure{\begin{tikzpicture}[baseline=-.6ex,x=1cm,y=1cm]
\draw (-.08,.42) .. controls (-.35,.25) and (.65,.47) .. (.58,-.24) -- (1.02,.65);
\end{tikzpicture}} sur $k$, et le fait que $A$ opère sur $M$ de façon que l'application diagonale soit $A$-linéaire signifie que $G$ opère sur $V$. On voit donc que si $u\handwritten{\unsure{\mapsto|\vdash}}\unsure{\deleted{U}u|\deleted{u}U}.e$ \inserted{de $A$ dans $\deleted{\unsure{N|V}}M$}\figure{\begin{tikzpicture}[baseline=-.6ex,x=1cm,y=1cm]
\draw (.06,-.24) .. controls (.02,.12) and (.34,.18) .. (.75,.2) -- (3.4,.2) .. controls (3.8,.22) and (3.8,.32) .. (3.76,.43);
\end{tikzpicture}} est injectif, i.e. si $g\handwritten{\mapsto}g.a$ ($a$ le \handwritten{``}centre\handwritten{''} de $V$) de $G$ dans $V$ est un monomorphisme, alors $M$ est libre sur $A$. En fait, on trouvera que $G$ opère librement sur $V$ i.e.
